% Supplementary Information for INFOGAIN
\documentclass[11pt]{article}
\usepackage[margin=1in]{geometry}
\usepackage{amsmath,amssymb,amsthm,mathtools}
\usepackage{graphicx,booktabs,caption}
\usepackage[expansion=false]{microtype}
\usepackage[hidelinks]{hyperref}
\usepackage[numbers,sort&compress]{natbib}

\captionsetup{font=small,labelfont=bf}
\graphicspath{{figures/}{../results/cohort_synthetic/mortality_30d/figures/}%
              {../results/simulation/figures/}}

\newtheorem{theorem}{Theorem}
\newtheorem{lemma}{Lemma}
\newtheorem{proposition}{Proposition}
\newtheorem{corollary}{Corollary}
\theoremstyle{definition}
\newtheorem{definition}{Definition}
\theoremstyle{remark}
\newtheorem{remark}{Remark}

\DeclareMathOperator{\KL}{KL}
\DeclareMathOperator{\TV}{TV}
\DeclareMathOperator{\NB}{NB}
\DeclareMathOperator{\Ber}{Ber}
\newcommand{\Vc}{\mathcal{V}}
\newcommand{\IV}{I_{\Vc}}
\newcommand{\HV}{H_{\Vc}}
\newcommand{\pvi}{\mathrm{pvi}}
\newcommand{\E}{\mathbb{E}}
\newcommand{\Prob}{\mathbb{P}}
\renewcommand{\thesection}{S\arabic{section}}
\renewcommand{\thetable}{S\arabic{table}}
\renewcommand{\thefigure}{S\arabic{figure}}


% Include a figure if the analysis that produces it has been run, otherwise leave
% a visible placeholder. This lets the manuscript compile from a partial results
% tree without ever silently dropping a panel.
\newcommand{\infogainfig}[2][\textwidth]{%
  \IfFileExists{#2.pdf}{\includegraphics[width=#1]{#2}}{%
    \IfFileExists{#2.png}{\includegraphics[width=#1]{#2}}{%
      \fbox{\parbox[c][0.26\textheight][c]{\dimexpr#1-2\fboxsep-2\fboxrule\relax}{%
        \centering\ttfamily\small
        [figure \detokenize{#2} not yet generated;
         run \detokenize{python scripts/run_all.py}]}}}}%
}

% Auto-generated by scripts/make_paper_numbers.py -- do not edit.
% Quantities from stages that have not run typeset as a visible marker
% rather than as a stale value or an undefined control sequence.
\providecommand{\todonum}[1]{\textbf{[?#1]}}

\newcommand{\thyAllHold}{yes}
\newcommand{\thyTrialsJS}{20000}
\newcommand{\thyIdentityRelErr}{1.0e-05}
\newcommand{\thyIdentityTrials}{150}
\newcommand{\thyPviErr}{1.7e-03}
\newcommand{\simRecovery}{82.9}
\newcommand{\simRecoveryCorrected}{96.0}
\newcommand{\simRegimeAccuracy}{100}
\newcommand{\simMaxN}{36{,}000}
\newcommand{\simThmOneHolds}{yes}
\newcommand{\simThmTwoHolds}{yes}
\newcommand{\simThmThreeRelErr}{3.8e-08}
\newcommand{\simRecoveryFloor}{0.01}
\newcommand{\simSubsetsScored}{11}
\newcommand{\simSubsetsBelowFloor}{0}
\newcommand{\simRecoveryFullPanel}{68.2}
\newcommand{\simMeanAbsErr}{0.0097}
\newcommand{\simRegimeMaxN}{50{,}000}
\newcommand{\simSynergyFirstN}{18{,}000}
\newcommand{\simRedundancyFirstN}{18{,}000}
\newcommand{\powBothFrac}{42}
\newcommand{\powEventsMin}{652}
\newcommand{\powEventsMax}{3569}
\newcommand{\powSynMin}{20}
\newcommand{\powSynMax}{82}
\newcommand{\powTreeSynMax}{72}
\newcommand{\restrictedShareMin}{52}
\newcommand{\restrictedShareMax}{98}
\newcommand{\powMidAN}{18{,}000}
\newcommand{\powMidAEvents}{1{,}391}
\newcommand{\powMidARecovery}{39}
\newcommand{\powMidBN}{50{,}000}
\newcommand{\powMidBEvents}{1{,}673}
\newcommand{\powMidBRecovery}{63}
\newcommand{\regimeSynSmall}{0}
\newcommand{\regimeSynSmallN}{6{,}000}
\newcommand{\regimeSynMid}{50}
\newcommand{\regimeSynMidN}{18{,}000}
\newcommand{\regimeSynLarge}{100}
\newcommand{\regimeSynLargeN}{50{,}000}
\newcommand{\acqLabs}{87}
\newcommand{\acqEcg}{67}
\newcommand{\acqCxr}{56}
\newcommand{\acqEcho}{23}
\newcommand{\thyThmOneInformative}{92.5}
\newcommand{\exGlobalNb}{0.146}
\newcommand{\exLocalNb}{0.061}
\newcommand{\tgtSignalRhoMin}{0.25}
\newcommand{\tgtSignalRhoMax}{0.50}
\newcommand{\tgtSignalCaptureMin}{65}
\newcommand{\tgtSignalCaptureMax}{74}
\newcommand{\tgtSignalChanceMin}{53}
\newcommand{\tgtSignalChanceMax}{54}
\newcommand{\tgtCeilingMax}{0.15}
\newcommand{\tgtN}{50{,}000}
\newcommand{\flipNPairs}{1}
\newcommand{\flipPairA}{laboratory panel}
\newcommand{\flipPairB}{electrocardiogram}
\newcommand{\flipSynEndpoint}{heart-failure readmission}
\newcommand{\flipRedEndpoint}{ICU transfer}
\newcommand{\flipSynBits}{0.0024}
\newcommand{\flipRedBits}{0.0004}
\newcommand{\ablSpeedup}{4}
\newcommand{\ablNCohorts}{2}
\newcommand{\switchNKeys}{691}
\newcommand{\switchMedianChange}{2.2}
\newcommand{\switchIncreased}{45}
\newcommand{\switchNSignFlips}{7}
\newcommand{\switchNMaterial}{0}
\newcommand{\switchNRegimeFlips}{4}
\newcommand{\permNPerm}{5{,}000}
\newcommand{\permFloor}{0.0002}
\newcommand{\decEchoNullPMin}{0.917}
\newcommand{\thyWitnessResidual}{6\times10^{-17}}
\newcommand{\certTightenMin}{3.2}
\newcommand{\certTightenMax}{9.9}
\newcommand{\shapleyEffResidual}{10^{-16}}
\newcommand{\redAvoidedMin}{8}
\newcommand{\redAvoidedMax}{75}
\newcommand{\redIcuTransferAvoided}{50}
\newcommand{\redIcuTransferSaved}{379}
\newcommand{\redIcuTransferCert}{0.0171}
\newcommand{\aurocIcuTransfer}{0.752}
\newcommand{\redAkiAvoided}{75}
\newcommand{\redAkiSaved}{419}
\newcommand{\redAkiCert}{0.0330}
\newcommand{\aurocAki}{0.762}
\newcommand{\redHfReadmissionAvoided}{8}
\newcommand{\redHfReadmissionSaved}{114}
\newcommand{\redHfReadmissionCert}{0.0164}
\newcommand{\aurocHfReadmission}{0.734}
\newcommand{\redMortalityAvoided}{29}
\newcommand{\redMortalitySaved}{352}
\newcommand{\redMortalityCert}{0.0160}
\newcommand{\aurocMortality}{0.798}
\newcommand{\cohortN}{50{,}000}
\newcommand{\cohortSubjects}{50{,}000}
\newcommand{\cohortSource}{synthetic}
\newcommand{\primaryPrevalence}{7.1}
\newcommand{\primaryBaselineBits}{0.0173}
\newcommand{\primaryFullBits}{0.0607}
\newcommand{\primaryBaselineAuroc}{0.665}
\newcommand{\primaryFullAuroc}{0.798}
\newcommand{\primaryNullGap}{0.00028}
\newcommand{\primaryIdentityRelErr}{3.9e-08}
\newcommand{\primaryRestrictedBits}{0.0329}
\newcommand{\primaryRestrictedShare}{94}
\newcommand{\primarySeedShare}{11}
\newcommand{\primaryMonoViolations}{13}
\newcommand{\primaryMonoWorst}{0.0009}
\newcommand{\redTestsAvoided}{29.0}
\newcommand{\redCostSaved}{351.64}
\newcommand{\redCostSavedPct}{77.3}
\newcommand{\redAuroc}{0.7961}
\newcommand{\redAurocAll}{0.7975}
\newcommand{\redCertifiedLoss}{0.01600}
\newcommand{\redInfoGap}{0.0015}
\newcommand{\gainMaxGini}{0.45}
\newcommand{\gainMinGini}{0.34}
\newcommand{\interNRedundant}{1}
\newcommand{\interNSynergistic}{1}
\newcommand{\interNPairs}{6}
\newcommand{\interMostRedundantPair}{labs--cxr}
\newcommand{\interMostRedundantBits}{0.0048}
\newcommand{\interRegimeFlips}{1}
\newcommand{\interNEndpoints}{4}
\newcommand{\tgtLabsRhoTruth}{0.25}
\newcommand{\tgtLabsRhoRetro}{0.075}
\newcommand{\tgtLabsCeiling}{0.150}
\newcommand{\tgtLabsTrueBits}{0.0382}
\newcommand{\tgtLabsCapture}{65}
\newcommand{\tgtLabsChance}{53}
\newcommand{\decLabsMarginal}{0.0265}
\newcommand{\decLabsConditional}{0.0231}
\newcommand{\decLabsRedundant}{0.0034}
\newcommand{\decLabsSynergistic}{0.0000}
\newcommand{\decLabsRedFrac}{13}
\newcommand{\gainLabsMean}{0.0243}
\newcommand{\gainLabsPNinetynine}{0.0875}
\newcommand{\gainLabsGini}{0.34}
\newcommand{\gainLabsConcentration}{4}
\newcommand{\gainLabsFracAbove}{90.3}
\newcommand{\tgtEcgRhoTruth}{0.11}
\newcommand{\tgtEcgRhoRetro}{-0.001}
\newcommand{\tgtEcgCeiling}{-0.008}
\newcommand{\tgtEcgTrueBits}{0.0002}
\newcommand{\tgtEcgCapture}{23}
\newcommand{\tgtEcgChance}{16}
\newcommand{\decEcgMarginal}{-0.0005}
\newcommand{\decEcgConditional}{0.0150}
\newcommand{\decEcgRedundant}{0.0000}
\newcommand{\decEcgSynergistic}{0.0155}
\newcommand{\decEcgRedFrac}{0}
\newcommand{\gainEcgMean}{0.0035}
\newcommand{\gainEcgPNinetynine}{0.0203}
\newcommand{\gainEcgGini}{0.45}
\newcommand{\gainEcgConcentration}{6}
\newcommand{\gainEcgFracAbove}{4.9}
\newcommand{\tgtCxrRhoTruth}{0.50}
\newcommand{\tgtCxrRhoRetro}{0.073}
\newcommand{\tgtCxrCeiling}{0.096}
\newcommand{\tgtCxrTrueBits}{0.0210}
\newcommand{\tgtCxrCapture}{74}
\newcommand{\tgtCxrChance}{54}
\newcommand{\decCxrMarginal}{0.0069}
\newcommand{\decCxrConditional}{0.0171}
\newcommand{\decCxrRedundant}{0.0000}
\newcommand{\decCxrSynergistic}{0.0103}
\newcommand{\decCxrRedFrac}{0}
\newcommand{\gainCxrMean}{0.0101}
\newcommand{\gainCxrPNinetynine}{0.0417}
\newcommand{\gainCxrGini}{0.44}
\newcommand{\gainCxrConcentration}{4}
\newcommand{\gainCxrFracAbove}{36.2}
\newcommand{\tgtEchoRhoTruth}{0.08}
\newcommand{\tgtEchoRhoRetro}{0.035}
\newcommand{\tgtEchoCeiling}{-0.001}
\newcommand{\tgtEchoTrueBits}{0.0002}
\newcommand{\tgtEchoCapture}{19}
\newcommand{\tgtEchoChance}{16}
\newcommand{\decEchoMarginal}{-0.0002}
\newcommand{\decEchoConditional}{-0.0005}
\newcommand{\decEchoRedundant}{0.0003}
\newcommand{\decEchoSynergistic}{0.0000}
\newcommand{\decEchoRedFrac}{176}
\newcommand{\gainEchoMean}{0.0030}
\newcommand{\gainEchoPNinetynine}{0.0146}
\newcommand{\gainEchoGini}{0.45}
\newcommand{\gainEchoConcentration}{5}
\newcommand{\gainEchoFracAbove}{3.9}
\newcommand{\tgtNotesRhoTruth}{\todonum{tgtNotesRhoTruth}}
\newcommand{\tgtNotesRhoRetro}{\todonum{tgtNotesRhoRetro}}
\newcommand{\tgtNotesCeiling}{\todonum{tgtNotesCeiling}}
\newcommand{\tgtNotesTrueBits}{\todonum{tgtNotesTrueBits}}
\newcommand{\tgtNotesCapture}{\todonum{tgtNotesCapture}}
\newcommand{\tgtNotesChance}{\todonum{tgtNotesChance}}
\newcommand{\decNotesMarginal}{\todonum{decNotesMarginal}}
\newcommand{\decNotesConditional}{\todonum{decNotesConditional}}
\newcommand{\decNotesRedundant}{\todonum{decNotesRedundant}}
\newcommand{\decNotesSynergistic}{\todonum{decNotesSynergistic}}
\newcommand{\decNotesRedFrac}{\todonum{decNotesRedFrac}}
\newcommand{\gainNotesMean}{\todonum{gainNotesMean}}
\newcommand{\gainNotesPNinetynine}{\todonum{gainNotesPNinetynine}}
\newcommand{\gainNotesGini}{\todonum{gainNotesGini}}
\newcommand{\gainNotesConcentration}{\todonum{gainNotesConcentration}}
\newcommand{\gainNotesFracAbove}{\todonum{gainNotesFracAbove}}
\newcommand{\ablAttnFmSyn}{58}
\newcommand{\ablAttnFmSynLo}{56}
\newcommand{\ablAttnFmSynHi}{59}
\newcommand{\ablAttnFmRegimes}{8/8}
\newcommand{\ablAttnFmSecs}{1{,}259}
\newcommand{\ablAttnNoFmSyn}{6}
\newcommand{\ablAttnNoFmSynLo}{6}
\newcommand{\ablAttnNoFmSynHi}{6}
\newcommand{\ablAttnNoFmRegimes}{4/8}
\newcommand{\ablAttnNoFmSecs}{870}
\newcommand{\ablConcatFmSyn}{57}
\newcommand{\ablConcatFmSynLo}{48}
\newcommand{\ablConcatFmSynHi}{66}
\newcommand{\ablConcatFmRegimes}{8/8}
\newcommand{\ablConcatFmSecs}{309}
\newcommand{\ablConcatNoFmSyn}{46}
\newcommand{\ablConcatNoFmSynLo}{39}
\newcommand{\ablConcatNoFmSynHi}{53}
\newcommand{\ablConcatNoFmRegimes}{8/8}
\newcommand{\ablConcatNoFmSecs}{334}
\newcommand{\decLabsMortalityMarg}{0.0265}
\newcommand{\decLabsMortalityCond}{0.0231}
\newcommand{\decLabsMortalityRed}{0.0034}
\newcommand{\decLabsMortalitySyn}{0.0000}
\newcommand{\decLabsMortalityP}{0.0002}
\newcommand{\decEcgMortalityMarg}{-0.0005}
\newcommand{\decEcgMortalityCond}{0.0150}
\newcommand{\decEcgMortalityRed}{0.0000}
\newcommand{\decEcgMortalitySyn}{0.0155}
\newcommand{\decEcgMortalityP}{0.0002}
\newcommand{\decCxrMortalityMarg}{0.0069}
\newcommand{\decCxrMortalityCond}{0.0171}
\newcommand{\decCxrMortalityRed}{0.0000}
\newcommand{\decCxrMortalitySyn}{0.0103}
\newcommand{\decCxrMortalityP}{0.0002}
\newcommand{\decEchoMortalityMarg}{-0.0002}
\newcommand{\decEchoMortalityCond}{-0.0005}
\newcommand{\decEchoMortalityRed}{0.0003}
\newcommand{\decEchoMortalitySyn}{0.0000}
\newcommand{\decEchoMortalityP}{0.9946}
\newcommand{\decNotesMortalityMarg}{\todonum{decNotesMortalityMarg}}
\newcommand{\decNotesMortalityCond}{\todonum{decNotesMortalityCond}}
\newcommand{\decNotesMortalityRed}{\todonum{decNotesMortalityRed}}
\newcommand{\decNotesMortalitySyn}{\todonum{decNotesMortalitySyn}}
\newcommand{\decNotesMortalityP}{\todonum{decNotesMortalityP}}
\newcommand{\decLabsHfReadmissionMarg}{0.0502}
\newcommand{\decLabsHfReadmissionCond}{0.0242}
\newcommand{\decLabsHfReadmissionRed}{0.0261}
\newcommand{\decLabsHfReadmissionSyn}{0.0000}
\newcommand{\decLabsHfReadmissionP}{0.0002}
\newcommand{\decEcgHfReadmissionMarg}{0.0103}
\newcommand{\decEcgHfReadmissionCond}{0.0156}
\newcommand{\decEcgHfReadmissionRed}{0.0000}
\newcommand{\decEcgHfReadmissionSyn}{0.0054}
\newcommand{\decEcgHfReadmissionP}{0.0002}
\newcommand{\decCxrHfReadmissionMarg}{0.0312}
\newcommand{\decCxrHfReadmissionCond}{0.0039}
\newcommand{\decCxrHfReadmissionRed}{0.0274}
\newcommand{\decCxrHfReadmissionSyn}{0.0000}
\newcommand{\decCxrHfReadmissionP}{0.0002}
\newcommand{\decEchoHfReadmissionMarg}{0.0090}
\newcommand{\decEchoHfReadmissionCond}{0.0126}
\newcommand{\decEchoHfReadmissionRed}{0.0000}
\newcommand{\decEchoHfReadmissionSyn}{0.0036}
\newcommand{\decEchoHfReadmissionP}{0.0002}
\newcommand{\decNotesHfReadmissionMarg}{\todonum{decNotesHfReadmissionMarg}}
\newcommand{\decNotesHfReadmissionCond}{\todonum{decNotesHfReadmissionCond}}
\newcommand{\decNotesHfReadmissionRed}{\todonum{decNotesHfReadmissionRed}}
\newcommand{\decNotesHfReadmissionSyn}{\todonum{decNotesHfReadmissionSyn}}
\newcommand{\decNotesHfReadmissionP}{\todonum{decNotesHfReadmissionP}}
\newcommand{\decLabsAkiMarg}{0.1215}
\newcommand{\decLabsAkiCond}{0.1143}
\newcommand{\decLabsAkiRed}{0.0072}
\newcommand{\decLabsAkiSyn}{0.0000}
\newcommand{\decLabsAkiP}{0.0002}
\newcommand{\decEcgAkiMarg}{-0.0004}
\newcommand{\decEcgAkiCond}{-0.0009}
\newcommand{\decEcgAkiRed}{0.0005}
\newcommand{\decEcgAkiSyn}{0.0000}
\newcommand{\decEcgAkiP}{0.9994}
\newcommand{\decCxrAkiMarg}{0.0091}
\newcommand{\decCxrAkiCond}{0.0014}
\newcommand{\decCxrAkiRed}{0.0078}
\newcommand{\decCxrAkiSyn}{0.0000}
\newcommand{\decCxrAkiP}{0.0036}
\newcommand{\decEchoAkiMarg}{-0.0006}
\newcommand{\decEchoAkiCond}{-0.0009}
\newcommand{\decEchoAkiRed}{0.0003}
\newcommand{\decEchoAkiSyn}{0.0000}
\newcommand{\decEchoAkiP}{1.0000}
\newcommand{\decNotesAkiMarg}{\todonum{decNotesAkiMarg}}
\newcommand{\decNotesAkiCond}{\todonum{decNotesAkiCond}}
\newcommand{\decNotesAkiRed}{\todonum{decNotesAkiRed}}
\newcommand{\decNotesAkiSyn}{\todonum{decNotesAkiSyn}}
\newcommand{\decNotesAkiP}{\todonum{decNotesAkiP}}
\newcommand{\decLabsIcuTransferMarg}{0.0149}
\newcommand{\decLabsIcuTransferCond}{0.0150}
\newcommand{\decLabsIcuTransferRed}{0.0000}
\newcommand{\decLabsIcuTransferSyn}{0.0001}
\newcommand{\decLabsIcuTransferP}{0.0002}
\newcommand{\decEcgIcuTransferMarg}{-0.0004}
\newcommand{\decEcgIcuTransferCond}{0.0007}
\newcommand{\decEcgIcuTransferRed}{0.0000}
\newcommand{\decEcgIcuTransferSyn}{0.0011}
\newcommand{\decEcgIcuTransferP}{0.0010}
\newcommand{\decCxrIcuTransferMarg}{0.0189}
\newcommand{\decCxrIcuTransferCond}{0.0208}
\newcommand{\decCxrIcuTransferRed}{0.0000}
\newcommand{\decCxrIcuTransferSyn}{0.0019}
\newcommand{\decCxrIcuTransferP}{0.0002}
\newcommand{\decEchoIcuTransferMarg}{-0.0001}
\newcommand{\decEchoIcuTransferCond}{-0.0002}
\newcommand{\decEchoIcuTransferRed}{0.0001}
\newcommand{\decEchoIcuTransferSyn}{0.0000}
\newcommand{\decEchoIcuTransferP}{0.9168}
\newcommand{\decNotesIcuTransferMarg}{\todonum{decNotesIcuTransferMarg}}
\newcommand{\decNotesIcuTransferCond}{\todonum{decNotesIcuTransferCond}}
\newcommand{\decNotesIcuTransferRed}{\todonum{decNotesIcuTransferRed}}
\newcommand{\decNotesIcuTransferSyn}{\todonum{decNotesIcuTransferSyn}}
\newcommand{\decNotesIcuTransferP}{\todonum{decNotesIcuTransferP}}


\title{\bfseries Supplementary Information\\[2mm]
INFOGAIN: decomposing the information value of multimodal diagnostic testing
at the level of the individual patient}
\date{}

\begin{document}
\maketitle
\tableofcontents
\clearpage

% ===================================================================== %
\section{Setting and standing assumptions}
\label{sec:setup}

Throughout, $Y\in\{0,1\}$ is the endpoint with prevalence
$\pi=\Prob(Y=1)\in(0,1)$, $M$ is a finite set of modalities, and
$X_S=(X_m)_{m\in S}$ for $S\subseteq M$. Logarithms are natural unless a
quantity is reported in bits, in which case it is divided by $\log 2$. We write
$H(\pi)=-\pi\log\pi-(1-\pi)\log(1-\pi)$ and $\eta_S(x)=\Prob(Y=1\mid X_S=x)$.

\begin{definition}[$\Vc$-usable information \citep{xu2020vinfo}]
For a family $\Vc$ of conditional distributions of $Y$ given an input (with the
convention that the input may be the empty tuple),
\[
\HV(Y)=\inf_{f\in\Vc}\E[-\log f[\varnothing](Y)],\qquad
\HV(Y\mid X_S)=\inf_{f\in\Vc}\E[-\log f[X_S](Y)],
\]
and $\IV(X_S\to Y)=\HV(Y)-\HV(Y\mid X_S)$.
\end{definition}

\begin{definition}[Masking closure]
$\Vc$ is \emph{closed under masking} if for every $f\in\Vc$ and every
$S\subseteq T\subseteq M$ there is $g\in\Vc$ with
$g[x_T]=f[x_S]$ for all $x$. The construction in
Methods~4 satisfies this by design: setting a presence flag to zero
substitutes a learned constant token for the modality's embedding, so the
resulting predictor is a member of the same parametric family.
\end{definition}

\begin{proposition}[Monotonicity]
\label{prop:mono}
If $\Vc$ is closed under masking then $S\subseteq T$ implies
$\HV(Y\mid X_T)\le \HV(Y\mid X_S)$ and hence
$\IV(X_S\to Y)\le \IV(X_T\to Y)$. Consequently the conditional gain
$U(m\mid S)=\IV(X_{S\cup m}\to Y)-\IV(X_S\to Y)$ is non-negative in population.
\end{proposition}

\begin{proof}
Every predictor available on $X_S$ is available on $X_T$ by masking the extra
coordinates, so the infimum defining $\HV(Y\mid X_T)$ is over a superset of the
achievable risks.
\end{proof}

\begin{remark}
Proposition~\ref{prop:mono} is a population statement. A finite sample and an
imperfectly optimised head can violate it by a few thousandths of a bit; we
report the number and size of such violations with every run and, alongside the
raw values, the Euclidean projection onto the monotone cone obtained by Dykstra's
cyclic projection over the covering relations of the subset lattice.
\end{remark}

\begin{remark}[Direction of the estimation error]
\label{rem:lowerbound}
Because $\HV(Y\mid X_S)$ is an infimum, any particular trained $\hat f_S$
satisfies $\E[-\log \hat f_S]\ge \HV(Y\mid X_S)$, so the plug-in
$\hat\IV$ \emph{under}-estimates $\IV$ up to sampling noise. Architectural or
optimisation weakness therefore costs tightness, never validity. Section~S5
makes this quantitative.
\end{remark}

% ===================================================================== %
\section{From information to discrimination: Lemmas and Theorem 1}
\label{sec:thm1}

\subsection{L1: the posterior maximises AUROC}

\begin{lemma}[Neyman--Pearson dominance]
\label{lem:np}
Let $A(s)$ denote the AUROC of a score $s$. Among all $X_T$-measurable scores,
$A$ is maximised by $\eta_T$. Write $A^*_T=A(\eta_T)$.
\end{lemma}

\begin{proof}
For any score $s$ and any level $\alpha$, the Neyman--Pearson lemma states that
the test rejecting on $\{\Lambda>c\}$, where $\Lambda=\mathrm{d}P_{X|1}/\mathrm{d}P_{X|0}$
and $c$ is chosen so the size is $\alpha$, has power at least that of any test
of size $\alpha$; in particular at least that of $\{s>t\}$ with
$\mathrm{FPR}(t)=\alpha$. Hence the ROC curve of $\Lambda$ dominates that of $s$
pointwise, and $A=\int_0^1\mathrm{TPR}\,\mathrm{d}\mathrm{FPR}$ inherits the
inequality. Finally $\eta_T=\pi\Lambda/(\pi\Lambda+1-\pi)$ is a strictly
increasing transform of $\Lambda$, and AUROC is invariant to strictly increasing
transforms.
\end{proof}

\subsection{L2: the Kolmogorov--Smirnov sandwich}

Let $F_1,F_0$ be the CDFs of the score under $Y=1$ and $Y=0$, and let
\[
D^+ \;=\; \sup_t\,[F_0(t)-F_1(t)] \;=\; \sup_t\,[\mathrm{TPR}(t)-\mathrm{FPR}(t)]\;\ge\;0
\]
be the one-sided Kolmogorov--Smirnov statistic (Youden's $J$ at its best
threshold; non-negative because the empty rejection region attains $0$).

\begin{lemma}[KS sandwich]
\label{lem:ks}
For any score, $2A-1\le 2D^+$. If in addition the ROC curve is concave, then
$2A-1\ge D^+$.
\end{lemma}

\begin{proof}
\emph{Upper bound.} With the mid-rank convention for ties,
$A=\int F_0\,\mathrm{d}F_1$ and $\int F_1\,\mathrm{d}F_1=\tfrac12$ exactly (also
in the presence of atoms). Since $F_0\le F_1+D^+$ pointwise,
$A\le \int F_1\,\mathrm{d}F_1+D^+=\tfrac12+D^+$.

\emph{Lower bound.} Let $(a,b)=(\mathrm{FPR}(t^*),\mathrm{TPR}(t^*))$ attain
$D^+=b-a$. A concave ROC through $(0,0)$ and $(1,1)$ lies above the chords
$(0,0)\to(a,b)$ and $(a,b)\to(1,1)$, whose combined area is
$\tfrac12 ab+\tfrac12(1-a)(1+b)=\tfrac12(1+b-a)$. Hence
$A\ge \tfrac12(1+D^+)$.
\end{proof}

\begin{remark}[Concavity is not decoration]
\label{rem:concave}
The lower bound genuinely fails for non-concave ROCs. Take positives all tied at
a score just below the median negative: then $A=\tfrac12$ while at the threshold
just below that common value $\mathrm{TPR}=1$, $\mathrm{FPR}=\tfrac12$, so
$D^+=\tfrac12>2A-1=0$. We therefore state and apply the lower bound for the ROC
convex hull \citep{provost2001robust}, which is achievable by a monotone
recalibration of the score and hence never exceeds $A^*$.
\end{remark}

\begin{lemma}
\label{lem:dtv}
For the posterior score, $D^+=\TV(P_{X|1},P_{X|0})$.
\end{lemma}

\begin{proof}
$\TV=\sup_B|P_1(B)-P_0(B)|$ is attained at $B=\{\Lambda>1\}$, a superlevel set of
$\eta$. The superlevel sets of $\eta$ are exactly the sets $\{\eta>t\}$, and
$P_1(\{\eta>t\})-P_0(\{\eta>t\})=\mathrm{TPR}(t)-\mathrm{FPR}(t)$; taking the
supremum over $t$ gives $D^+=\TV$.
\end{proof}

\subsection{L3: information against total variation}

Two elementary identities drive both directions. Since $\E[\eta]=\pi$,
\begin{equation}
\label{eq:tv-eta}
\TV(P_{X|1},P_{X|0})=\tfrac12\!\int\!|p_1-p_0| = \frac{\E|\eta-\pi|}{2\pi(1-\pi)},
\qquad
I(X;Y)=H(\pi)-\E[H(\eta)] = \E\big[\KL(\Ber(\eta)\,\|\,\Ber(\pi))\big].
\end{equation}
The first follows from $p_1=\eta p/\pi$ and $p_0=(1-\eta)p/(1-\pi)$, giving
$p_1-p_0=p(\eta-\pi)/(\pi(1-\pi))$; the second by expanding the $\KL$ and using
$\E[\eta]=\pi$.

\begin{lemma}[L3]
\label{lem:jstv}
With $\TV=\TV(P_{X|1},P_{X|0})$ and $I=I(X;Y)$ in nats,
\[
8\pi^2(1-\pi)^2\,\TV^2 \;\le\; I \;\le\; H(\pi)\,\TV .
\]
\end{lemma}

\begin{proof}
\emph{Upper bound.} Put $\varphi(\eta)=H(\pi)-H(\eta)$, so $I=\E[\varphi(\eta)]$
by \eqref{eq:tv-eta}. The binary entropy $H$ is concave with $H(0)=H(1)=0$, so it
lies above each of its chords: on $[0,\pi]$, $H(\eta)\ge (\eta/\pi)H(\pi)$, and
on $[\pi,1]$, $H(\eta)\ge H(\pi)(1-\eta)/(1-\pi)$. Hence
\[
\varphi(\eta)\;\le\; H(\pi)\left[\frac{(\pi-\eta)_+}{\pi}+\frac{(\eta-\pi)_+}{1-\pi}\right].
\]
Because $\E[\eta]=\pi$ we have $\E[(\pi-\eta)_+]=\E[(\eta-\pi)_+]=\tfrac12\E|\eta-\pi|=:D$,
so taking expectations,
\[
I \;\le\; H(\pi)\,D\left[\frac1\pi+\frac1{1-\pi}\right]
  \;=\; \frac{H(\pi)\,D}{\pi(1-\pi)} \;=\; H(\pi)\,\TV,
\]
the last step by \eqref{eq:tv-eta}, since $\TV=2D/(2\pi(1-\pi))=D/(\pi(1-\pi))$.

\emph{Lower bound.} Pinsker's inequality for Bernoulli laws gives
$\KL(\Ber(\eta)\|\Ber(\pi))\ge 2(\eta-\pi)^2$, so by \eqref{eq:tv-eta} and
Jensen,
\[
I \;\ge\; 2\,\E[(\eta-\pi)^2]\;\ge\;2\big(\E|\eta-\pi|\big)^2
  \;=\;2\big(2\pi(1-\pi)\TV\big)^2 \;=\; 8\pi^2(1-\pi)^2\TV^2 . \qedhere
\]
\end{proof}

Both inequalities are tight in the relevant limits: the upper bound with equality
when $\eta\in\{0,1\}$ (disjoint class-conditionals), the lower bound
asymptotically when $|\eta-\pi|$ is constant. In $2\times10^{4}$ Monte-Carlo
trials over Dirichlet class-conditionals with random $\pi$ neither was violated;
the worst slack was $5\times10^{-14}$, i.e.\ machine precision
(\texttt{infogain.theory.lemmas.verify\_js\_tv}).

\begin{remark}[A sharper constant we do not use]
The stronger inequality $I\ge 2\pi(1-\pi)\TV^2$ --- larger by a factor
$1/(4\pi(1-\pi))\ge1$ --- also held in every trial. Reducing it to two-point laws
by Bauer's maximum principle (the functional
$(\E|\eta-\pi|)^2/(2\pi(1-\pi))-\E[\KL]$ is convex in the law, and the extreme
points of the set of laws on $[0,1]$ with mean $\pi$ are supported on at most two
points) leaves a two-parameter inequality we have not proved analytically. We
therefore state and use only the Pinsker version. Nothing depends on the
difference: the sharper constant would only tighten the \emph{upper} end of the
AUROC bracket, and a looser upper bound is still an upper bound.
\end{remark}

\begin{corollary}[Information bracket for achievable AUROC]
\label{cor:bracket}
$\dfrac12+\dfrac{I}{2H(\pi)}\;\le\;A^*\;\le\;\dfrac12+\dfrac{\sqrt I}{2\sqrt2\,\pi(1-\pi)}$,
both sides clipped to $[\tfrac12,1]$.
\end{corollary}

\begin{proof}
Combine Lemmas~\ref{lem:np}--\ref{lem:jstv}: the lower bound from
$2A^*-1\ge D^+=\TV\ge I/H(\pi)$, the upper from
$2A^*-1\le 2D^+=2\TV\le 2\sqrt{I/(8\pi^2(1-\pi)^2)}$.
\end{proof}

\subsection{Theorem 1: synergy lower-bounds the achievable gain}

\begin{theorem}[Stratified achievable-AUROC gain]
\label{thm:s-auroc}
Let $\{B_k\}_{k=1}^K$ be any partition of the range of $\eta_S$ into intervals,
let $w_k$ be the fraction of (case, control) pairs with both members in $B_k$,
and let $A_k(\cdot)$ denote AUROC computed within $B_k$. Then
\[
A^*_{S\cup m}\;\ge\;A(\eta_S)\;+\;\sum_{k=1}^{K} w_k\big[A_k(\eta_{S\cup m})-A_k(\eta_S)\big].
\]
In population $A(\eta_S)=A^*_S$ by Lemma~\ref{lem:np}, so the sum lower-bounds
the achievable gain $A^*_{S\cup m}-A^*_S$.
\end{theorem}

\begin{proof}
Index the bins in increasing order of $\eta_S$ and define the lexicographic score
\[
s_\lambda(x) \;=\; k(x) \;+\; \lambda\,\eta_{S\cup m}(x),
\qquad \lambda \in \Big(0,\ \big[\sup \eta_{S\cup m}-\inf \eta_{S\cup m}\big]^{-1}\Big),
\]
where $k(x)$ is the index of the bin containing $\eta_S(x)$. It is
$X_{S\cup m}$-measurable, so $A^*_{S\cup m}\ge A(s_\lambda)$ by
Lemma~\ref{lem:np}.

Decompose $A$ over ordered pairs (one case, one control). For a pair in
different bins, the choice of $\lambda$ makes the bin index decide the
comparison, and since bins are intervals ordered by $\eta_S$ the decision agrees
with the one $\eta_S$ makes; such pairs therefore contribute identically to
$A(s_\lambda)$ and to $A(\eta_S)$. For a pair inside bin $k$ the bin indices tie
and $\eta_{S\cup m}$ decides. Writing $w_k$ for the share of pairs inside bin $k$
and $1-\sum_k w_k$ for the share across bins,
\[
A(s_\lambda) \;=\; \underbrace{\Big[A(\eta_S)-\sum_k w_k A_k(\eta_S)\Big]}_{\text{across-bin part, unchanged}}
\;+\;\sum_k w_k A_k(\eta_{S\cup m}),
\]
which rearranges to the claim.
\end{proof}

\begin{corollary}[Information form]
\label{cor:thm1-info}
Applying Corollary~\ref{cor:bracket} inside each bin, with $\pi_k$ the within-bin
prevalence and $U_k=\E[U(m\mid S)\mid B_k]$ the within-bin conditional usable
information,
\[
A^*_{S\cup m}-A^*_S \;\ge\; \sum_k w_k\left[\frac{U_k}{2H(\pi_k)}-\big(A_k(\eta_S)-\tfrac12\big)\right]_+ .
\]
As the partition refines, $A_k(\eta_S)\to\tfrac12$ and the bound is driven purely
by the within-bin conditional information, i.e.\ by the synergy term.
\end{corollary}

\begin{remark}[Why a stratified statement is necessary]
AUROC is not a proper scoring rule, so no unconditional inequality of the form
``$\Delta I\ge\varepsilon \Rightarrow \Delta A \ge g(\varepsilon)>0$'' can hold:
a model can gain information purely by recalibration while leaving every ranking
--- and hence AUROC --- untouched. Theorem~\ref{thm:s-auroc} sidesteps this by
bounding what is \emph{achievable} rather than what a particular score achieves,
and by localising to the strata inside which the ranking must change.
\end{remark}

\paragraph{Numerical verification.} The proof's content is the displayed
\emph{identity} between the certified sum and the AUROC gain of the witness
score $s_\lambda$, and that is what we check: across synergy-dominated
simulations the two agreed to $5\times10^{-17}$, i.e.\ machine precision. The
weaker implied inequality (bound $\le$ gain demonstrably achievable on the same
rows) held with zero violations, and the bound was strictly informative
($>0.005$ AUROC) in over 90\% of trials. On exact simulator posteriors it held
for every modality (\simThmOneHolds).

\paragraph{Choosing the partition.} The theorem holds for any \emph{fixed}
partition, so maximising the bound over partitions on the same data it then
certifies is a selection effect. By default the number of bins is chosen on one
random half of the cohort and the bound reported on the other, which makes the
partition independent of the data it certifies at some cost in power.

% ===================================================================== %
\section{Theorem 3: the information--net-benefit identity}
\label{sec:thm3}

Let $\eta_1=\eta_S$ and $\eta_2=\eta_{S\cup m}$, and put
$V_j(t)=\E[(\eta_j-t)_+]$, $\NB_j(t)=V_j(t)/(1-t)$.

\begin{lemma}[Net benefit of the Bayes rule]
\label{lem:nb}
The rule ``treat iff $\eta>t$'' has net benefit
$\NB(t)=\E[(\eta-t)_+]/(1-t)$, where net benefit is
$\mathrm{TP}/n-(\mathrm{FP}/n)\cdot t/(1-t)$ \citep{vickers2006decision}.
\end{lemma}

\begin{proof}
With $w=t/(1-t)$, $\E[\mathrm{TP}/n]=\E[\mathbf 1\{\eta>t\}\eta]$ and
$\E[\mathrm{FP}/n]=\E[\mathbf 1\{\eta>t\}(1-\eta)]$, so
$\NB(t)=\E[\mathbf 1\{\eta>t\}(\eta(1+w)-w)]=(1+w)\E[(\eta-t)_+]$
because $w/(1+w)=t$; and $1+w=1/(1-t)$.
\end{proof}

\begin{lemma}[Schervish representation for the binary log-loss]
\label{lem:scherv}
Let $G(p)=p\log p+(1-p)\log(1-p)$ and $\nu(t)=G''(t)=1/(t(1-t))$. For
$p,q\in(0,1)$,
\[
\KL(\Ber(p)\|\Ber(q)) \;=\; \int_0^1\Big[(p-t)_+-(q-t)_+-\mathbf 1\{q>t\}(p-q)\Big]\nu(t)\,\mathrm{d}t .
\]
\end{lemma}

\begin{proof}
The bracket vanishes for $t<\min(p,q)$ (it equals $(p-t)-(q-t)-(p-q)=0$) and for
$t>\max(p,q)$ (all three terms vanish), so the integrand is supported on the
interval between $p$ and $q$ and the integral converges despite
$\int_0\nu=\infty$. On $q<t<p$ the bracket is $p-t$; on $p<t<q$ it is $t-p$; in
both cases $|p-t|$. Hence the right-hand side equals
$\int_{\min}^{\max}|p-t|\nu(t)\,\mathrm{d}t$, which is the Bregman divergence of
$G$ evaluated at $(p,q)$ by the integral form of Taylor's theorem
$G(p)-G(q)-G'(q)(p-q)=\int_q^p (p-t)G''(t)\,\mathrm{d}t$; and the Bregman
divergence of the negative binary entropy is exactly
$\KL(\Ber(p)\|\Ber(q))$ \citep{schervish1989general}.
\end{proof}

\begin{theorem}[Information--net-benefit identity]
\label{thm:s-identity}
If $\eta_1=\E[\eta_2\mid X_S]$ (the nesting condition), then in nats
\[
\Delta I \;:=\; \E\,\KL(\Ber(\eta_2)\|\Ber(\eta_1))
\;=\;\int_0^1\big[V_2(t)-V_1(t)\big]\frac{\mathrm{d}t}{t(1-t)}
\;=\;\int_0^1\big[\NB_2(t)-\NB_1(t)\big]\frac{\mathrm{d}t}{t}.
\]
\end{theorem}

\begin{proof}
Take expectations in Lemma~\ref{lem:scherv} with $p=\eta_2$, $q=\eta_1$. The
third term has conditional expectation
$\E[\mathbf 1\{\eta_1>t\}(\eta_2-\eta_1)\mid X_S]=\mathbf 1\{\eta_1>t\}\,
\E[\eta_2-\eta_1\mid X_S]=0$ by the nesting condition, since
$\mathbf 1\{\eta_1>t\}$ is $X_S$-measurable. Fubini is licensed because the
integrand is non-negative after this cancellation (by
Proposition~\ref{prop:mono}, or directly by Jensen: $V_2\ge V_1$ since
$(\cdot-t)_+$ is convex). The first equality follows; the second is
Lemma~\ref{lem:nb}.
\end{proof}

\begin{remark}[Convergence at the endpoints]
Near $t=0$, $V_j(t)=\pi-\E[\min(\eta_j,t)]$, so
$V_2(t)-V_1(t)=\E[\min(\eta_1,t)]-\E[\min(\eta_2,t)]$, which is $O(t)$; the
integrand $\Delta V/(t(1-t))$ is therefore bounded there. Near $t=1$,
$\Delta V(t)\le V_2(t)\le (1-t)$, so the integrand is again bounded.
\end{remark}

\begin{definition}[Clinically restricted usable information]
For a threshold window $[a,b]\subset(0,1)$,
$I^{[a,b]}:=\int_a^b \Delta\NB(t)\,\mathrm{d}t/t$.
\end{definition}

Since $\Delta\NB\ge0$, $I^{[a,b]}\le\Delta I$, with the deficit being bits
purchased at thresholds outside clinical practice. Theorem~\ref{thm:s-identity}
also gives immediately
$\sup_{t\in[a,b]}\Delta\NB(t)\ \ge\ I^{[a,b]}/\log(b/a)$: information gained in
the window must show up as net benefit somewhere in it.

\paragraph{Numerical verification.} Over \thyIdentityTrials\ random nested
posterior pairs of $3\times10^4$ patients each, the maximum relative discrepancy
between the two sides was \thyIdentityRelErr; on the fitted cohort models it was
\primaryIdentityRelErr\ (Fig.~7 of the main text).

% ===================================================================== %
\section{Theorem 2: certified safe omission}
\label{sec:thm2}

\begin{theorem}[Global bound]
\label{thm:s-omit-global}
Suppose omitting modality $m$ forgoes $\Delta I\le\varepsilon$ nats. Then for
every $t\in(0,1)$, $\Delta\NB(t)\le\sqrt{\varepsilon/2}\,/\,(1-t)$.
\end{theorem}

\begin{proof}
Write $\Delta V=V_2-V_1\ge0$. Its derivative is
$\Delta V'(t)=\Prob(\eta_1>t)-\Prob(\eta_2>t)$, so $|\Delta V'|\le1$: $\Delta V$
is $1$-Lipschitz. Since $\nu(t)=1/(t(1-t))\ge4$ on $(0,1)$,
Theorem~\ref{thm:s-identity} gives
$\int_0^1\Delta V(t)\,\mathrm{d}t\le \varepsilon/4$. A non-negative $1$-Lipschitz
function with supremum $h$ attained at $t_0$ dominates the tent
$(h-|t-t_0|)_+$, whose integral is at least $h^2/2$ (exactly $h^2$ for an
interior maximum, $h^2/2$ at an endpoint). Hence
$h^2/2\le\varepsilon/4$, i.e.\ $h\le\sqrt{\varepsilon/2}$, and
$\Delta\NB(t)=\Delta V(t)/(1-t)$.
\end{proof}

\begin{theorem}[Threshold-localised bound]
\label{thm:s-omit-local}
Fix $t_0\in(0,1)$ and let $v^\star(\varepsilon,t_0)$ be the unique solution of
\[
\int_{t_0-v}^{t_0+v}\big(v-|t-t_0|\big)\,\frac{\mathrm{d}t}{t(1-t)} \;=\; \varepsilon
\]
(intersected with $(0,1)$). Then $\Delta V(t_0)\le v^\star$ and
$\Delta\NB(t_0)\le v^\star/(1-t_0)$.
\end{theorem}

\begin{proof}
Identical to Theorem~\ref{thm:s-omit-global} but retaining the true weight rather
than the bound $\nu\ge4$: if $\Delta V(t_0)=v$ then, by non-negativity and
$1$-Lipschitzness, $\Delta V(t)\ge (v-|t-t_0|)_+$ everywhere, so
$\varepsilon=\int\Delta V\nu\ge\int (v-|t-t_0|)_+\nu$. The left-hand side of the
displayed equation is continuous and strictly increasing in $v$, giving
uniqueness.
\end{proof}

Because $\nu$ blows up at small $t$, Theorem~\ref{thm:s-omit-local} is
substantially tighter exactly where triage decisions live: at $t_0=0.10$ and
$\varepsilon$ corresponding to $0.05$ bits, the global bound gives
$\Delta\NB\le0.146$ and the localised bound $\Delta\NB\le0.061$.

\begin{corollary}[Cohort-level omission]
The per-patient bound is concave in $\varepsilon$, so by Jensen the
cohort-average net-benefit loss from omitting $m$ for a set of patients is at
most $v^\star(\bar\varepsilon,t_0)/(1-t_0)$ with $\bar\varepsilon$ the average
forgone information over those patients; patients for whom the test is kept
contribute zero.
\end{corollary}

\begin{corollary}[Study planning]
Inverting $v^\star$ gives the minimum information gain a test must deliver before
it could move the decision curve at threshold $t_0$ by a clinically visible
amount --- a power calculation in bits.
\end{corollary}

Both bounds dominated the realised loss in all $300$ Monte-Carlo trials and, on
exact simulator posteriors, for every modality and threshold (\simThmTwoHolds).

% ===================================================================== %
\section{Properties of the context-relative decomposition}
\label{sec:decomp-props}

\begin{proposition}[Exactness and disjointness]
With $I_m$, $U(m\mid S)$, $R(m\mid S)$, $\mathrm{Syn}(m\mid S)$ as in
Section~2.1 of the main text,
$U(m\mid S)=I_m-R(m\mid S)+\mathrm{Syn}(m\mid S)$ identically,
$R,\mathrm{Syn}\ge0$, and $R\cdot\mathrm{Syn}=0$.
\end{proposition}

\begin{proof}
Immediate from $x=(x)_+-(-x)_+$ applied to $x=U-I_m$.
\end{proof}

\begin{remark}[Relation to partial information decomposition]
The classical PID lattice \citep{williams2010pid} requires a redundancy
functional that its axioms do not determine, and the leading candidates
\citep{bertschinger2014uniqueinfo,griffith2014synergy} disagree on real data. We
avoid the choice rather than make it: by conditioning on the context actually in
hand, $R$ and $\mathrm{Syn}$ become differences of separately estimable
quantities. The price is that they are properties of $(m,S)$ rather than of the
modalities alone --- which is precisely the clinical content, since the value of
a test is not a property of the test.
\end{remark}

\begin{proposition}[Shapley attribution]
Let $v(T)=\IV(X_{S_0\cup T}\to Y)-\IV(X_{S_0}\to Y)$ for $T$ a subset of the
orderable modalities, and
$\varphi_m=\sum_{T\subseteq M_o\setminus m}\frac{|T|!\,(|M_o|-|T|-1)!}{|M_o|!}
[v(T\cup m)-v(T)]$. Then $\varphi$ is the unique attribution satisfying
efficiency, symmetry, null-player and additivity \citep{shapley1953value}, and
$\sum_m\varphi_m=v(M_o)$.
\end{proposition}

Because $v$ is itself a sample mean of pointwise values, $\varphi_m$ is a mean of
\emph{per-patient} Shapley contributions, which yields a patient-level
attribution and a bootstrap interval at no extra cost. Efficiency is verified
numerically at each run; the residual is below $10^{-15}$ bits.

% ===================================================================== %
\section{Finite-sample behaviour of the estimator}
\label{sec:finite}

Let $\pvi_i(S)=\log_2 f_S(y_i\mid x_i^S)-\log_2 f_\varnothing(y_i)$, clipped to
$[-C,C]$ with $C=8$ bits, evaluated out of fold.

\begin{theorem}[Two-sided bracket]
\label{thm:s-bracket}
Let $\hat I=\frac1n\sum_i\pvi_i(S)$ and let $\hat V$ be its sample variance.
With probability at least $1-\delta$,
\[
\IV(X_S\to Y)\;\ge\;\hat I-\varepsilon_{\mathrm{EB}},\qquad
\IV(X_S\to Y)\;\le\;\hat I+\varepsilon_{\mathrm{EB}}+2L_\ell\mathfrak R_n(\Vc)+\gamma_{\mathrm{opt}},
\]
where
$\varepsilon_{\mathrm{EB}}=\sqrt{2\hat V\log(4/\delta)/n}+\tfrac{7\cdot 2C\log(4/\delta)}{3(n-1)}$
is the empirical-Bernstein deviation \citep{maurer2009empirical},
$\mathfrak R_n(\Vc)$ the Rademacher complexity of the head class, $L_\ell$ the
Lipschitz constant of the loss in the logit on the clipped region, and
$\gamma_{\mathrm{opt}}$ the optimisation gap.
\end{theorem}

\begin{proof}
The lower bound needs no capacity control: $\hat f_S\in\Vc$ witnesses
$\HV(Y\mid X_S)\le\E[-\log \hat f_S]$ (Remark~\ref{rem:lowerbound}), and
$\varepsilon_{\mathrm{EB}}$ bounds the deviation of the out-of-fold sample mean of
a variable in an interval of width $2C$. The upper bound adds the uniform
convergence term needed to pass from the trained predictor's risk to the family
infimum, plus the residual optimisation gap.
\end{proof}

We instantiate $\mathfrak R_n$ by the norm-based bound of
\citet{golowich2018size}, $\mathfrak R_n\le B\big(\prod_{l}\|W_l\|_F\big)
\sqrt{2L\log 2}/\sqrt n$, computed from the trained weights. It is loose in
absolute terms, as all norm-based bounds are, but it is computable, monotone in
the quantities weight decay controls, and it makes the upper side a checkable
statement rather than an asymptotic gesture.

\paragraph{Bias.} Two biases act in opposite directions. Evaluation bias
(upward) is removed by construction through cross-fitting
\citep{chernozhukov2018dml}. Learning bias (downward) is the dominant term: it
is the gap between what a finite training set taught the head and the family
infimum. We measure it by refitting at several training fractions and
extrapolating $\hat I(n)=I_\infty-a n^{-\alpha}$ with $\alpha\in[0.15,2]$, in the
tradition of \citet{strong1998entropy} and \citet{panzeri2007correcting} but with
a free exponent, because forcing $\alpha=1$ badly over-extrapolates the
interaction terms that carry synergy.

\paragraph{Variance.} Total estimator variance is decomposed into patient
sampling and seed-to-seed variability; the latter accounted for
\primarySeedShare\% of the total on the primary endpoint and is added to every
interval. Differences (conditional gains, marginals, Shapley values) are computed
pairwise per patient before averaging, which cancels the shared null term and the
patient-level variance that dominates each component separately.

% ===================================================================== %
\section{Identification of the per-patient counterfactual}
\label{sec:ident}

$\Delta_i(m\mid S)$ asks what a test \emph{would have} shown for a patient. For
patients who received the test this is an interpolation; for those who did not it
is a counterfactual, identified under two conditions.

\begin{proposition}[Identification]
Let $O_m\in\{0,1\}$ indicate that $m$ was acquired. If
(i) $X_m \perp O_m \mid X_S$ (missing at random given the observed context
\citep{rubin1976inference}) and
(ii) $0<\Prob(O_m=1\mid X_S=x)$ for almost every $x$ in the support (overlap),
then $q(x_m\mid x^S)=p(x_m\mid x^S,O_m=1)$ and $\Delta_i$ is identified from the
observed data.
\end{proposition}

Condition (ii) is checkable and we check it: a cross-validated propensity model
for $O_m$ given the baseline context is fitted per modality, and the fraction of
patients outside the region of common support is reported (Table~S5). Patients
there are flagged rather than scored --- for them the counterfactual is an
extrapolation, and reporting a number would be worse than reporting none.

Condition (i) is not checkable, because it concerns the unobserved result. We
therefore bound its consequences by simulation: the generator's
\texttt{mnar\_strength} parameter makes ordering depend on latent factors that
the model cannot see but the ordering clinician can, and the whole analysis is
repeated. The bias has a consistent direction --- unmeasured indication inflates
the apparent value of selectively-ordered tests --- which makes the reduction
claim conservative: a policy fitted under MNAR orders more tests than are needed,
not fewer.

\paragraph{Consistency check.} By the tower property,
$\E_i[\Delta_i(m\mid S)]$ and the retrospective realised gain
$\E_i[\pvi_i(S\cup m)-\pvi_i(S)]$ estimate the same population quantity by
completely different routes --- the first integrating over hypothetical results
without using labels, the second using realised results and labels. Their
agreement (Table~S4) is a joint test of the conditional sampler, the masked
family and the cross-fitting, and is reported for every modality.


% ===================================================================== %
\section{Raw-signal representations}
\label{sec:raw}

The headline analysis represents the electrocardiogram by its interval and axis
measurements plus report keyword flags, and the chest radiograph by its CheXpert
label vector. Both are deliberate choices, and both cost usable information.

The case for them is auditability and access. A cardiologist can read a row of
the ECG block and say whether the decomposition is plausible; nobody can do that
with the 48th coordinate of a learned embedding. And the tabular path needs
about 1\,GB of metadata rather than the roughly 660\,GB of waveforms and pixels,
which is the difference between a study another group can reproduce and one they
cannot.

The case against them is that they discard information the raw signal contains.
Because $\IV$ is defined relative to the predictor family, substituting a
stronger encoder \emph{raises} it --- that is the intended semantics, not a
confound, and it is why the family must be reported alongside the number.
\texttt{infogain.encoders.waveform} implements a 1-D residual encoder over
10\,s $\times$ 12 leads at 100\,Hz, and \texttt{infogain.encoders.image} a
compact 2-D residual encoder over $224^2$ radiographs; both emit an embedding of
the same width as the tabular blocks and are masked identically, so the family
stays closed under masking and every theorem applies unchanged.
\texttt{embeddings\_from\_npz} additionally accepts externally computed
per-study embeddings, which is the hook for a pretrained thoracic model.

The comparison --- how many of the bits the tabular representation forgoes ---
requires the full waveform and image downloads and is reported once the
credentialed extraction has been run. We expect the direction to be unambiguous
(more information) and the decomposition's \emph{qualitative} structure to be
stable, since redundancy between the laboratory panel and the radiograph is a
statement about shared pathophysiology rather than about encoder capacity; but
that expectation is a hypothesis this analysis will test, not a result.

% ===================================================================== %
\section{Supplementary tables}

\begin{table}[htbp]\centering\small
\caption{Numerical certification of every theoretical claim
(\texttt{python -m infogain.theory.verify}). ``Worst slack'' is the smallest
margin by which the inequality held across all trials; negative values of order
$10^{-13}$ reflect floating-point rounding at an equality case. All claims
pass: \thyAllHold.}
\label{tab:s-theory}
\begin{tabular}{lrrrl}
\toprule
\textbf{claim} & \textbf{trials} & \textbf{worst slack} & \textbf{violations} & \textbf{holds} \\
\midrule
I \textless{}= H(pi) * TV & 20000 & 5.05e-14 & 0 & yes \\
I \textgreater{}= 2 pi (1-pi) TV\textasciicircum{}2 & 20000 & 3.25e-14 & 0 & yes \\
2A\_hull-1 \textgreater{}= D+ & 2000 & -2.08e-17 & 0 & yes \\
2A-1 \textless{}= 2 D+ & 2000 & 1.02e-02 & 0 & yes \\
Thm3 identity: dI = int dNB dt/t & 150 & -1.01e-05 & 0 & yes \\
Thm2 global bound dominates dNB & 300 & 9.86e-03 & 0 & yes \\
Thm2 local bound dominates dNB & 300 & 4.15e-03 & 0 & yes \\
Thm1 bound \textless{}= achievable gain & 200 & 1.36e-03 & 0 & yes \\
Thm1 bound informative (\textgreater{}0.005 AUROC) & 200 & 9.15e-01 & 0 & yes \\
PVI recovers I(X;Y) (truth=0.17312, est=0.17146) & 400000 & -1.66e-03 & 0 & yes \\
\bottomrule
\end{tabular}

\end{table}

\begin{table}[htbp]\centering\small
\caption{Simulation study: estimated against exact information by cohort size,
with the learning-curve-corrected estimate. Recovery approaches 100\% from
below, as a computationally bounded quantity must.}
\label{tab:s-recovery}
\begin{tabular}{rlrrrrrr}
\toprule
\textbf{n} & \textbf{subset} & \textbf{true bits} & \textbf{est bits} & \textbf{CI lo} & \textbf{CI hi} & \textbf{recovered \textbackslash{}\%} & \textbf{corrected} \\
\midrule
4000 & demo+vitals & 0.0083 & 0.0027 & -0.0070 & 0.0123 & 32.7342 & 0.0144 \\
4000 & cxr+demo+ecg+echo+labs+vitals & 0.0769 & 0.0213 & 0.0043 & 0.0378 & 27.7398 & 0.0520 \\
4000 & demo+labs+vitals & 0.0491 & 0.0231 & 0.0130 & 0.0335 & 47.1035 & 0.0431 \\
4000 & demo+ecg+vitals & 0.0090 & 0.0005 & -0.0090 & 0.0091 & 5.9845 & 0.0132 \\
4000 & cxr+demo+vitals & 0.0226 & 0.0087 & -0.0014 & 0.0189 & 38.6418 & 0.0242 \\
4000 & demo+echo+vitals & 0.0085 & 0.0013 & -0.0103 & 0.0127 & 15.0333 & 0.0140 \\
4000 & cxr+demo+ecg+echo+vitals & 0.0459 & 0.0051 & -0.0074 & 0.0171 & 11.0418 & 0.0215 \\
4000 & cxr+demo+echo+labs+vitals & 0.0525 & 0.0236 & 0.0068 & 0.0404 & 45.0382 & 0.0521 \\
4000 & demo+ecg+echo+labs+vitals & 0.0491 & 0.0191 & 0.0082 & 0.0298 & 38.9042 & 0.0401 \\
4000 & cxr+demo+ecg+labs+vitals & 0.0770 & 0.0224 & 0.0066 & 0.0378 & 29.0660 & 0.0553 \\
4000 & demo+ecg+labs+vitals & 0.0491 & 0.0209 & 0.0110 & 0.0310 & 42.6030 & 0.0419 \\
12000 & demo+vitals & 0.0074 & 0.0133 & 0.0045 & 0.0223 & 180.1819 & 0.0144 \\
12000 & cxr+demo+ecg+echo+labs+vitals & 0.0713 & 0.0339 & 0.0255 & 0.0427 & 47.5460 & 0.0520 \\
12000 & demo+labs+vitals & 0.0429 & 0.0336 & 0.0273 & 0.0402 & 78.4591 & 0.0431 \\
12000 & demo+ecg+vitals & 0.0081 & 0.0128 & 0.0050 & 0.0207 & 157.5048 & 0.0132 \\
12000 & cxr+demo+vitals & 0.0197 & 0.0187 & 0.0098 & 0.0278 & 94.7713 & 0.0242 \\
12000 & demo+echo+vitals & 0.0076 & 0.0132 & 0.0063 & 0.0203 & 174.6777 & 0.0140 \\
12000 & cxr+demo+ecg+echo+vitals & 0.0439 & 0.0169 & 0.0079 & 0.0263 & 38.6078 & 0.0215 \\
12000 & cxr+demo+echo+labs+vitals & 0.0461 & 0.0355 & 0.0286 & 0.0430 & 76.9900 & 0.0521 \\
12000 & demo+ecg+echo+labs+vitals & 0.0430 & 0.0319 & 0.0260 & 0.0380 & 74.1682 & 0.0401 \\
12000 & cxr+demo+ecg+labs+vitals & 0.0713 & 0.0343 & 0.0265 & 0.0423 & 48.0493 & 0.0553 \\
12000 & demo+ecg+labs+vitals & 0.0430 & 0.0322 & 0.0263 & 0.0382 & 74.9073 & 0.0419 \\
36000 & demo+vitals & 0.0150 & 0.0141 & 0.0102 & 0.0179 & 93.8918 & 0.0144 \\
36000 & cxr+demo+ecg+echo+labs+vitals & 0.0749 & 0.0413 & 0.0346 & 0.0480 & 55.1519 & 0.0520 \\
36000 & demo+labs+vitals & 0.0477 & 0.0386 & 0.0324 & 0.0448 & 80.9906 & 0.0431 \\
36000 & demo+ecg+vitals & 0.0151 & 0.0126 & 0.0085 & 0.0165 & 83.6328 & 0.0132 \\
36000 & cxr+demo+vitals & 0.0273 & 0.0222 & 0.0174 & 0.0271 & 81.3913 & 0.0242 \\
36000 & demo+echo+vitals & 0.0151 & 0.0136 & 0.0093 & 0.0178 & 90.1061 & 0.0140 \\
36000 & cxr+demo+ecg+echo+vitals & 0.0493 & 0.0202 & 0.0151 & 0.0254 & 41.0314 & 0.0215 \\
36000 & cxr+demo+echo+labs+vitals & 0.0517 & 0.0424 & 0.0352 & 0.0496 & 82.1183 & 0.0521 \\
36000 & demo+ecg+echo+labs+vitals & 0.0477 & 0.0369 & 0.0308 & 0.0429 & 77.2958 & 0.0401 \\
36000 & cxr+demo+ecg+labs+vitals & 0.0749 & 0.0418 & 0.0354 & 0.0482 & 55.8721 & 0.0553 \\
36000 & demo+ecg+labs+vitals & 0.0477 & 0.0374 & 0.0318 & 0.0431 & 78.4801 & 0.0419 \\
\bottomrule
\end{tabular}

\end{table}

\begin{table}[htbp]\centering\small
\caption{Regime recovery: true against estimated decomposition regime
(redundant / synergistic / independent) per modality and cohort size. Accuracy
at the largest cohort is \simRegimeAccuracy\%. Redundancy, a main effect, is
recovered first; synergy, an interaction across two noisy read-outs, last.}
\label{tab:s-regime}
\begin{tabular}{lrrrrllc}
\toprule
\textbf{modality} & \textbf{true \$I\_m\$} & \textbf{est \$I\_m\$} & \textbf{true \$U\$} & \textbf{est \$U\$} & \textbf{true regime} & \textbf{est regime} & \textbf{correct} \\
\midrule
labs & 0.0340 & 0.0273 & 0.0272 & 0.0209 & redundant & redundant & True \\
ecg & 0.0002 & -0.0010 & 0.0241 & -0.0018 & synergistic & independent & False \\
cxr & 0.0127 & 0.0118 & 0.0286 & 0.0058 & synergistic & redundant & False \\
echo & 0.0001 & -0.0011 & 0.0000 & -0.0009 & independent & independent & True \\
\bottomrule
\end{tabular}

\end{table}

\begin{table}[htbp]\centering\small
\caption{Smallest cohort at which each modality's regime is called correctly,
against the size of the effect being detected. This is the planning number: it
says what it costs to be able to tell a redundant test from a synergistic one at
a given effect size.}
\label{tab:s-firstcorrect}
\begin{tabular}{lrrr}
\toprule
\textbf{modality} & \textbf{true regime} & \textbf{true effect bits} & \textbf{first correct n} \\
\midrule
ecg & synergistic & 0.0220 & 18000 \\
cxr & synergistic & 0.0147 & 50000 \\
labs & redundant & 0.0060 & 18000 \\
echo & independent & 0.0001 & 6000 \\
\bottomrule
\end{tabular}

\end{table}

\begin{table}[htbp]\centering\small
\caption{Prospective against retrospective per-patient gain. The two are
computed by different routes and must agree in cohort mean by the tower
property.}
\label{tab:s-consistency}
\begin{tabular}{lrrrr}
\toprule
\textbf{modality} & \textbf{prospective mean bits} & \textbf{retrospective mean bits} & \textbf{abs difference} & \textbf{spearman} \\
\midrule
labs & 0.0243 & 0.0262 & 0.0018 & 0.1097 \\
ecg & 0.0035 & -0.0006 & 0.0041 & -0.0009 \\
cxr & 0.0101 & 0.0065 & 0.0037 & 0.0770 \\
echo & 0.0030 & -0.0002 & 0.0033 & 0.0305 \\
\bottomrule
\end{tabular}

\end{table}

\begin{table}[htbp]\centering\small
\caption{Overlap diagnostics for the counterfactual. A propensity model that
predicts ordering almost perfectly is a warning, not a success: it means there
is a region of patient space where the test is never ordered.}
\label{tab:s-overlap}
\begin{tabular}{lrrrrrr}
\toprule
\textbf{modality} & \textbf{order rate} & \textbf{propensity auroc} & \textbf{frac off support} & \textbf{propensity p01} & \textbf{propensity p99} & \textbf{ess ratio} \\
\midrule
labs & 0.8714 & 0.7224 & 0.0260 & 0.5300 & 0.9856 & 0.9876 \\
ecg & 0.6704 & 0.7565 & 0.0029 & 0.1578 & 0.9687 & 0.9167 \\
cxr & 0.5560 & 0.7844 & 0.0034 & 0.0625 & 0.9643 & 0.8380 \\
echo & 0.2320 & 0.8083 & 0.0554 & 0.0072 & 0.8431 & 0.5600 \\
\bottomrule
\end{tabular}

\end{table}

\begin{table}[htbp]\centering\small
\caption{Full subset lattice of $\Vc$-usable information, with monotonicity
diagnostics.}
\label{tab:s-lattice}
\begin{tabular}{lrrrr}
\toprule
\textbf{subset} & \textbf{$|S|$} & \textbf{bits} & \textbf{CI lo} & \textbf{CI hi} \\
\midrule
EMPTY & 0 & -0.0003 & -0.0010 & 0.0004 \\
cxr & 1 & 0.0067 & 0.0053 & 0.0080 \\
ecg & 1 & -0.0002 & -0.0013 & 0.0008 \\
echo & 1 & 0.0000 & -0.0007 & 0.0007 \\
labs & 1 & 0.0242 & 0.0213 & 0.0270 \\
demo+vitals & 2 & 0.0173 & 0.0151 & 0.0196 \\
cxr+demo+vitals & 3 & 0.0242 & 0.0216 & 0.0268 \\
demo+ecg+vitals & 3 & 0.0168 & 0.0140 & 0.0196 \\
demo+echo+vitals & 3 & 0.0171 & 0.0149 & 0.0194 \\
demo+labs+vitals & 3 & 0.0438 & 0.0403 & 0.0472 \\
cxr+demo+ecg+vitals & 4 & 0.0382 & 0.0341 & 0.0423 \\
cxr+demo+echo+vitals & 4 & 0.0238 & 0.0213 & 0.0265 \\
cxr+demo+labs+vitals & 4 & 0.0458 & 0.0422 & 0.0495 \\
demo+ecg+echo+vitals & 4 & 0.0164 & 0.0137 & 0.0192 \\
demo+ecg+labs+vitals & 4 & 0.0438 & 0.0399 & 0.0477 \\
demo+echo+labs+vitals & 4 & 0.0438 & 0.0403 & 0.0472 \\
cxr+demo+ecg+echo+vitals & 5 & 0.0376 & 0.0337 & 0.0414 \\
cxr+demo+ecg+labs+vitals & 5 & 0.0612 & 0.0564 & 0.0659 \\
cxr+demo+echo+labs+vitals & 5 & 0.0457 & 0.0420 & 0.0494 \\
demo+ecg+echo+labs+vitals & 5 & 0.0436 & 0.0396 & 0.0476 \\
cxr+demo+ecg+echo+labs+vitals & 6 & 0.0607 & 0.0562 & 0.0653 \\
\bottomrule
\end{tabular}

\end{table}

% ===================================================================== %
\clearpage
\section{Supplementary figures}

\begin{figure}[htbp]\centering
\infogainfig[\textwidth]{figS1_recovery}
\caption{\textbf{Recovering known information.} Left: estimated against exact
information at the largest simulated cohort, before and after learning-curve
correction. Right: fraction recovered against cohort size, over subsets whose
true information exceeds 0.01 bits. The estimator approaches the truth from
below, as a computationally bounded quantity must.}
\label{fig:s-recovery}
\end{figure}

\begin{figure}[htbp]\centering
\infogainfig[\textwidth]{figS2_regime}
\caption{\textbf{Recovering the decomposition's sign.} Left: share of
modalities assigned the correct regime against cohort size, overall and split by
the true regime. Redundancy is a main effect and is recovered first; synergy is
an interaction between latent factors in two different noisy read-outs and is
recovered last. Right: the same information as a planning number -- effect size
against the cohort needed to detect it.}
\label{fig:s-regime}
\end{figure}

\bibliographystyle{unsrtnat}
\bibliography{refs}
\end{document}
